\chapter{实机调试部署技巧}
\label{chap:deployment}

\begin{introduction}
  \item 学习目标：选择现场调试入口；用 SSH、远程桌面和主机名连接机器人；配置有线直连与 Tailscale；同步明确版本并核验运行；使用会话、服务、日志与 KVM 工具定位故障。
  \item 建议学时：约 2--3 天，包含两机连接、心跳项目与故障排查；真实设备与长时间运行按队内计划另行验收。
  \item 本卷边界：以 Ubuntu 22.04 机载电脑与常见桌面客户端为例，讲部署和维护基本功；相机、CUDA、ROS 2 及控制链路见卷二、卷三。工具资料按 2026 年 10 月 31 日核对，软件界面和硬件型号应以实际版本为准。
\end{introduction}

\begingroup
\setlength{\emergencystretch}{3em}
\clubpenalty=10000
\widowpenalty=10000

\section{代码走出自己的电脑}
\label{sec:deploy-physical}

前面几章里，你已经能写代码、编译工程，还能请 AI 帮忙排查问题。接下来，学长把机器人推过来：“放到车上跑一下吧。”这时候你才发现，\textbf{在自己的电脑上能运行，只是交付的起点}。机器人真正使用的是机载电脑（onboard computer）：程序要在它的处理器和系统上运行，读取它的相机与串口，并在重新上电后按预期工作。你电脑上的文件夹不会因为两台机器都装了 Ubuntu 就自动出现在车上，x86 电脑编译的程序也不能直接拿给 ARM 板运行。

我们把准备目标环境、传送代码与配置、在目标机构建或安装、启动并核验结果这一过程称为\textbf{部署（deployment）}。本章先解决一个最实际的问题：怎样操作机载电脑？等连接稳定了，再讨论如何把程序放上去、让它持续运行，并在故障时找回控制权。

\subsubsection{便携屏和键鼠：最直接的入口}

最传统的办法，是给机载电脑接上一块便携显示器，再插 USB 键盘与鼠标。图~\ref{fig:deploy-monitor} 是一块带 HDMI 输入的小屏幕；对机载电脑来说，它就是普通外接显示器。你可以看到开机画面、登录桌面，像操作自己的电脑一样打开终端，网络还没配好时尤其有用。

\begin{figure}[htbp]
  \centering
  \includegraphics[width=0.58\textwidth]{deployment/portable-monitor.jpg}
  \caption{便携 HDMI 显示器的正反面。先确认视频输入与供电接口，再准备键鼠；图源：Waveshare 官方产品页。}
  \label{fig:deploy-monitor}
\end{figure}

不过，机器人通常不在桌上。你可能一边蹲在地上扶屏幕，一边伸手够键盘，旁边还拖着供电线和视频线；换个位置调车，整套“移动办公桌”又得搬一次。便携屏非常适合首次装机和故障救援，日常开发则可以让自己的笔记本承担键盘与屏幕，把终端或桌面通过连接传回来。接下来的 SSH 和远程桌面，就是为了让你少蹲一会儿。

\begin{example}[首次接入机载电脑]
前提是机载电脑已有可启动的系统，机器人处于队内确认的安全调试状态，不会因为程序误启动产生运动。
先核对电脑的视频\emph{输出}接口与显示器的\emph{输入}接口：HDMI 对 HDMI；只有 DP 输出时，使用方向正确的 DP 转 HDMI 适配器。按显示器铭牌准备电源，再接视频线，最后把有线键鼠插到机载电脑的 USB 口。USB-C 的外形不保证支持视频输出，也不保证能为屏幕供电，应查看两端说明书。

开机后在屏幕上选择对应输入源，确认能看到启动或登录画面。进入桌面后打开终端，运行下面两条只读命令，记录用户名与主机名；后文以用户 \texttt{rm}、主机名 \texttt{rm-robot} 为例，实际使用时换成自己的。
\begin{tbash}
whoami
hostname
\end{tbash}
如果屏幕显示“无信号”，先检查输入源、线材方向与供电，再看主机是否实际启动；此时还没进入 SSH，修改网络配置通常帮不上忙。
\end{example}

\begin{exercise}\label{ex:deploy-physical}
给队内一台机载电脑画出便携屏、屏幕电源与 USB 键鼠的接线图，标明输入和输出方向。接入后记录实际用户名、主机名以及视频接口；没有设备时先按型号手册完成接线方案，现场验证留到本章验收。
\end{exercise}

实物接口可对照\href{https://www.waveshare.com/13.3inch-FHD-Monitor.htm}{Waveshare 官方产品页}查看。教材选图用于认识便携屏的外观；HDMI 输入与供电接口的位置应放大产品接口图核对，不要求购买同一型号。

\section{先把终端搬过来：SSH}
\label{sec:deploy-ssh}

便携屏能让你看到机器人上的桌面，但如果只是改参数、编译、看日志，每次搬一套外设实在有点累。安全外壳协议（Secure Shell，SSH）可以让笔记本打开机载电脑的终端：你在自己的键盘上输入，命令却在机器人上执行，结果再传回来。它传输的是终端数据，不需要不断发送整张桌面，通常很适合远程开发。

这里先假定两台电脑已接入同一台路由器，机器人当前地址为 \texttt{192.168.50.20}，普通用户名为 \texttt{rm}。这两个值都要换成自己的；如何找到地址、用主机名代替地址以及没有路由器时怎么办，稍后在第~\ref{sec:deploy-lan}、\ref{sec:deploy-wired}~节动手。\textbf{能上网不代表两台设备之间能互通}，访客网络可能禁止客户端相互访问。

\subsubsection{装在谁上面，命令又在哪边执行}

SSH 分为客户端与服务端：笔记本运行客户端发起连接，机载电脑运行服务端接收连接。Ubuntu 上 \texttt{openssh-client} 提供前者，\texttt{openssh-server} 提供后者。Windows 10/11 通常可在 PowerShell 或终端使用 \texttt{ssh}；若找不到命令，在“可选功能”安装 OpenSSH 客户端。机载电脑装的是 Ubuntu Server 还是 Desktop，不影响使用 SSH。

\begin{example}[从笔记本进入机载电脑]
先借助便携屏，在\textbf{机器人本机终端}安装并核对服务：
\begin{tbash}
sudo apt update
sudo apt install openssh-server
sudo systemctl enable --now ssh
systemctl is-active ssh
hostname
whoami
ip -br address
sudo ssh-keygen -lf /etc/ssh/ssh_host_ed25519_key.pub
\end{tbash}
\texttt{is-active} 应输出 \texttt{active}；记下实际用户名、当前接口的 IPv4 地址，以及最后一行的 \texttt{SHA256:...} 主机指纹。若系统使用防火墙，应由维护者允许调试网络到 SSH 端口，不必为了排错关闭整个防火墙。

然后在\textbf{笔记本终端}执行，基本形式是 \texttt{ssh 用户名@地址}：
\begin{tbash}
ssh rm@192.168.50.20
hostname
whoami
pwd
exit
\end{tbash}
第一次连接会问是否信任主机密钥。把显示的指纹与刚才在机器人本机看到的指纹核对，一致再输入 \texttt{yes}，然后输入机器人上 \texttt{rm} 用户的密码；输入密码时没有星号回显是正常现象。连接后，\texttt{hostname} 应是机器人而非笔记本的名字，\texttt{pwd} 通常是远端家目录。\texttt{exit} 退出远端终端，回到笔记本。
\end{example}

如果提示 \texttt{Connection timed out}，先查地址、链路、客户端隔离与防火墙；\texttt{Connection refused} 通常表示目标拒绝该端口连接，应检查 SSH 服务和端口；\texttt{Permission denied} 则优先核对用户名、密码或密钥。用 \texttt{ssh -v rm@地址} 查看连接过程。终端提示符可自定义，容易看花眼时，主动运行 \texttt{hostname} 与 \texttt{pwd} 确认操作对象。

\begin{note}
主机密钥验证的是“你连到哪台机器”。如果以前连接过，后来出现 \texttt{REMOTE HOST IDENTIFICATION HAS CHANGED}，先核对是否重装系统、地址被另一台机器占用或连接目标错误；确认新指纹后再按文档更新记录。不要养成禁用主机密钥检查的习惯。改 SSH 配置前也应保留当前连接与本地救援入口，先运行 \texttt{sudo sshd -t} 检查配置。
\end{note}

\subsubsection{密钥和连接别名：少输一点，也更好管理}

用户密钥回答“你是谁”，与刚才用于核验服务器身份的主机密钥职责不同。私钥留在笔记本，公钥可以交给机器人；私钥密码用于保护本地私钥，也不等于机器人登录密码。下面以 Ubuntu 笔记本为例；Windows 客户端也可生成密钥，分发公钥的方法见 OpenSSH 文档。

\begin{example}[给自己的调试连接建档]
在笔记本上生成专用密钥；若该名字已存在，选新名字，不覆盖旧密钥。按提示设置私钥密码：
\begin{tbash}
ssh-keygen -t ed25519 -f ~/.ssh/id_ed25519_rm -C "rm-debug-laptop"
ssh-copy-id -i ~/.ssh/id_ed25519_rm.pub rm@192.168.50.20
ssh -i ~/.ssh/id_ed25519_rm rm@192.168.50.20
\end{tbash}
确认可登录后，退出连接，把下面一段加入\textbf{笔记本}已有的 \path{~/.ssh/config}，不要覆盖其他配置：
\begin{tbash}
Host rm-robot
    HostName 192.168.50.20
    User rm
    IdentityFile ~/.ssh/id_ed25519_rm
    IdentitiesOnly yes
    ServerAliveInterval 15
    ServerAliveCountMax 3
\end{tbash}
之后运行 \texttt{ssh rm-robot} 即可。\texttt{Host} 是本机别名，\texttt{HostName} 指定实际目标。以后可将目标改为 \path{rm-robot.local} 或 Tailscale 名称。\texttt{IdentityFile} 指明使用哪份私钥；保活参数帮助发现连接中断，不会自动重启远端程序。用 \texttt{ssh -G rm-robot} 可展开当前实际采用的客户端配置。
\end{example}

不要把私钥提交到 Git，也不要让全队共用一份个人私钥。每人保留自己的身份，离队或笔记本丢失时，维护者便能单独撤销相应公钥。公钥登录先在另一个终端验证成功，再考虑是否调整密码登录策略，避免把自己锁在门外。

\begin{exercise}\label{ex:deploy-ssh}
使用自己的两台设备完成一次 SSH 登录，记录客户端与服务端的主机名、用户名和主机指纹核验方式。再添加一个连接别名，以密钥登录，说明公钥、私钥、主机密钥分别保存在哪里、负责验证谁。只有一台机器时可先在虚拟机练习，并明确标注不是实机验收。
\end{exercise}

\begin{conclusion}
SSH 让你远程使用终端；它需要网络可达、远端服务启动及身份验证成功。IP 地址或可解析的名字负责找机器，主机指纹负责确认机器，用户名与用户密钥负责登录。断开 SSH 后程序是否继续运行，还取决于会话与进程的启动方式，见第~\ref{sec:deploy-ops}~节。
\end{conclusion}

安装与服务检查可查 \href{https://ubuntu.com/server/docs/how-to/security/openssh-server/}{Ubuntu OpenSSH 指南}；密钥和客户端配置参数分别见 \href{https://man.openbsd.org/ssh-keygen}{ssh-keygen 手册}、\href{https://man.openbsd.org/ssh_config}{ssh\_config 手册}。

\subsubsection{想保留熟悉的编辑器：VS Code Remote -- SSH}

SSH 不要求你只能用终端编辑器。第 1 章用过的 VS Code 可安装 Microsoft 的 \texttt{Remote - SSH} 扩展，在笔记本显示界面，在机载电脑读写工程、运行终端与调试器。这属于远程开发：它不会直接把机器人的整个桌面搬过来，也不会自动把所有本地文件同步过去。

\begin{example}[确认编辑与编译发生在哪台电脑]
先在普通终端验证 \texttt{ssh rm-robot} 成功，再在\textbf{笔记本}的 VS Code 扩展页安装 \texttt{Remote - SSH}。打开命令面板，选择 \texttt{Remote-SSH: Connect to Host...}，选择配置好的 \texttt{rm-robot}。首次连接会在远端安装 VS Code Server，需要可用下载通路及满足扩展要求的系统与资源；离线场地应提前准备。

连接后检查状态栏显示远程主机，再用 \texttt{File > Open Folder} 打开\textbf{远端}工程目录。在新建的集成终端运行 \texttt{hostname}、\texttt{pwd} 和 \texttt{uname -m}，应看到机器人的主机名、目录和体系结构。打开一个练习文本，加入一行并保存；在独立 SSH 终端用 \texttt{cat} 读取相同的远端路径，确认修改落在远端。C++、CMake 等开发扩展按提示安装到远端，编译器与依赖也需要在远端准备。
\end{example}

\begin{exercise}\label{ex:deploy-remote-editor}
用远程窗口打开一个练习工程，从远端集成终端构建或运行本章心跳示例。把 \texttt{hostname} 与当前路径记入日志，再指出本地窗口与远程窗口各修改哪个文件系统。若安装失败，保留扩展日志，区分 SSH 连接失败与远端 Server 安装失败。
\end{exercise}

新版本对远端 Linux 的库版本、架构与下载条件可能有不同要求，安装前核对 \href{https://code.visualstudio.com/docs/remote/ssh}{Remote -- SSH 官方教程}和其中的系统要求；这也适用于资源有限的 ARM 机载电脑。

\section{需要看画面时：远程桌面}
\label{sec:deploy-desktop}

SSH 很适合编译、查看日志和启动程序，但有些操作确实需要窗口：例如打开相机厂商的配置工具、查看图形调试器，或演示桌面上的设置。这时可以使用\textbf{远程桌面（remote desktop）}：机载电脑上的服务把桌面画面发给笔记本，笔记本再把键盘与鼠标操作传回去。它依赖机器人上的系统、桌面会话与服务正常工作；后面介绍的硬件 KVM 走的是另一条路径。

\subsubsection{Ubuntu 的 RDP 与笔记本上的客户端}

\textbf{远程桌面协议（Remote Desktop Protocol，RDP）}是一种传输桌面与输入的协议。\textbf{Remmina} 是 Linux 上的客户端，Windows 则可以用自带的“远程桌面连接” \texttt{mstsc}。只在笔记本上安装客户端还不够，\emph{被连接的机载电脑}也要启用对应的服务。

本卷以 Ubuntu 22.04 的 GNOME 桌面为例：在“设置 → 共享”中开启共享，再进入“远程桌面”。图~\ref{fig:deploy-gnome-rdp} 是 GNOME 42 官方发布说明中的界面，重点看远程桌面、远程控制和认证三组设置；图中系统与主机名是官方示例，不是你的连接地址。Ubuntu 后续版本可能把入口移到“设置 → 系统 → 远程桌面”，并区分“桌面共享”与“远程登录”，不要把新版本的功能直接套回 22.04。

\begin{figure}[htbp]
  \centering
  \includegraphics[width=0.62\textwidth]{deployment/gnome42-rdp.png}
  \caption{GNOME 42 的 RDP 设置官方截图。Remote Desktop 允许连接，Remote Control 决定能否操作；连接凭据在 Authentication 中设置。}
  \label{fig:deploy-gnome-rdp}
\end{figure}

\begin{example}[从笔记本操作 Ubuntu 22.04 桌面]
这个示例要求双方已有可达的 IP 地址；暂时不知道地址时，先完成第~\ref{sec:deploy-lan}~节的网络核对。机载电脑先接便携屏并登录 GNOME，保持当前桌面处于解锁、唤醒状态。
\begin{enumerate}
  \item 在机载电脑“设置 → 共享 → 远程桌面”中开启 Remote Desktop 和 Remote Control，记下界面给出的地址、用户名与密码。这里使用的是\emph{远程桌面设置中的凭据}，不要想当然地填入 Linux 登录密码；查看“验证加密（Verify Encryption）”中的指纹。
  \item Ubuntu 笔记本安装并打开 Remmina：
\begin{tbash}
sudo apt update
sudo apt install remmina remmina-plugin-rdp
remmina
\end{tbash}
  新建连接，协议选 RDP，服务器填机器人 IP（例如 \texttt{192.168.50.20}）；用户名、密码填上一步的 RDP 凭据，保存并连接。已有其他网络时先确认该示例地址属于当前机器人，不能直接照抄。
  \item Windows 笔记本可按 \texttt{Win+R}，输入 \texttt{mstsc}，在“计算机”中填同一个 IP，再输入 RDP 凭据。若界面给出非默认端口，应填 \texttt{地址:端口}；本例服务通常使用 TCP 3389。
  \item 首次连接先将客户端提示的证书指纹与机载电脑上的记录比较。进入远程桌面后打开终端，运行 \texttt{hostname}，并移动一个窗口；同时确认机载屏幕上能看到对应变化，说明连到的是预期会话，且有控制权限。
\end{enumerate}
如果双方启用了防火墙，由管理员仅为已核对的调试网段或 Tailscale 接口开放所需服务。不要为了试连直接关闭整机防火墙，也不要把 RDP 端口转发到公网。完成实验后可回到共享设置关闭远程桌面。
\end{example}

\begin{note}
22.04 的上述入口主要共享已经登录的桌面。退出登录、锁屏、显示器休眠或拔掉显示器后是否仍能工作，要在本机实际测试；不能从“这次能连”推导出“下次无屏上电也能连”。纯命令行系统、不同桌面、Wayland/X11 会话和显卡驱动也可能需要不同方案。先用便携屏确认会话状态；需要可靠的无头启动时，优先保留 SSH，再按对应版本文档配置桌面服务。
\end{note}

\subsubsection{RustDesk：跨平台的另一种入口}

\textbf{RustDesk} 也常用于跨平台远程操作。与直接填写 RDP 地址不同，它常用设备 ID 寻找另一端，并通过公共服务或队内自建服务协助连接。图~\ref{fig:deploy-rustdesk} 中左侧是本机 ID 与一次性密码，右侧是要连接的设备 ID；这些字段属于官方演示，不要使用图中的号码连接陌生设备。

\begin{figure}[htbp]
  \centering
  \includegraphics[width=0.82\textwidth]{deployment/rustdesk-client.png}
  \caption{RustDesk 官方客户端截图。左侧信息给对方连接本机使用，右侧输入目标设备 ID；底部状态提示当前是否已就绪。}
  \label{fig:deploy-rustdesk}
\end{figure}

\begin{example}[用 RustDesk 建立一次受控连接]
在双方电脑上从\href{https://rustdesk.com/docs/en/client/}{官方客户端文档}进入下载页，选择与系统及处理器架构对应的安装包。Ubuntu 用其 \texttt{.deb} 包；在下载目录运行 \texttt{sudo apt install ./实际文件名.deb}，这里的文件名必须换成你刚下载的包，不能同时随意安装多个版本。

先在机载电脑本地打开 RustDesk，确认底部显示就绪；在笔记本上输入\emph{机载电脑的} ID，点击 Connect。由机载电脑一侧确认本次请求，或使用该次显示的一次性密码完成认证。连接后打开终端核对 \texttt{hostname}，检查鼠标与键盘是否能操作；然后主动断开，并在机载电脑确认连接已经结束。无人值守的固定密码与队内自建服务属于后续配置，先把这次有人确认的连接做通。
\end{example}

RustDesk 的 Linux 官方说明仍把 Wayland 支持列为实验性，远程访问登录界面要求 X11。若需要切换，可在本地退出当前会话、于登录页选择“Ubuntu on Xorg”后重新测试；登录管理器与无人值守设置另按\href{https://rustdesk.com/docs/en/client/linux/}{官方 Linux 文档}核对，不能只切换一次桌面就承诺重启后的登录界面可用。无头设备还要确认系统实际提供了可捕获的显示输出。它也不能替代网络配置：依赖公共 ID 服务时要能访问相应服务，脱离互联网后的使用必须提前配置并验证。

\begin{exercise}\label{ex:deploy-desktop}
选择 RDP 或 RustDesk，完成一次“笔记本连入 → 核对主机名 → 操作窗口 → 主动断开”的记录，再分别测试锁屏与退出登录后的连接表现。说明你使用的是共享现有桌面还是另一种会话，记录系统、桌面与工具版本。没有第二台电脑时，先列出所需条件和验证方法，不能把预期写成实测。
\end{exercise}

\begin{conclusion}
RDP 是协议，Remmina 与 \texttt{mstsc} 是客户端；RustDesk 是另一套远程访问工具。选择时先问“目标机是否已有可用桌面、我要连接现有会话还是登录新会话、断网和重启后怎么办”。图形连接失败时先保留 SSH 查服务与日志；能移动远程鼠标，只能证明远程输入和显示链路通了，真实相机帧率与控制延迟仍要在设备端测量。
\end{conclusion}

Ubuntu 22.04 的 RDP 支持见\href{https://documentation.ubuntu.com/release-notes/22.04/}{该版本发布说明}；设置与 Remmina 连接步骤可对照\href{https://ubuntu.com/project/docs/maintainers/test-plans/desktop-team/RemoteDesktop/}{Ubuntu 官方远程桌面测试计划}，新版本的入口差异则查看\href{https://help.ubuntu.com/stable/ubuntu-help/sharing-desktop.html.zh-CN}{当前桌面共享帮助}。

\section{局域网：先找到机器人，再记住它的名字}
\label{sec:deploy-lan}

远程连接首先得回答一个朴素的问题：机器人在哪里？这里问的是网络地址。把笔记本和机载电脑接到队内路由器后，路由器通常用动态主机配置协议（Dynamic Host Configuration Protocol，DHCP）为它们分配地址。例如，笔记本是 \texttt{192.168.50.10/24}，机器人是 \texttt{192.168.50.20/24}；\texttt{/24} 对应子网掩码 \texttt{255.255.255.0}，这个例子里二者属于同一子网。\textbf{这些只是示例地址，动手时以实际输出为准}，不要把书上的地址直接写进已有网络。

在机载电脑本地终端运行 \texttt{ip -br address}，查看处于 \texttt{UP} 状态的网口及其地址；\texttt{lo} 上的 \texttt{127.0.0.1} 只指向本机，不能拿给笔记本连接。\texttt{ip route} 查看路由，含 \texttt{default via} 的一行表示默认出口。只在同一子网内调试不要求有互联网，但两台设备必须能够互相通信：校园网、访客 Wi-Fi 或路由器的客户端隔离，可能让“连着同一个 Wi-Fi”仍然无法互访。

\begin{example}[按地址检查局域网连通性]
先在机器人本地查看：
\begin{tbash}
ip -br address
ip route
\end{tbash}
假设确认有线地址是 \texttt{192.168.50.20/24}，再到笔记本的 Linux 终端运行：
\begin{tbash}
ip -br address
ip route get 192.168.50.20
ping -c 3 192.168.50.20
ssh rm@192.168.50.20
\end{tbash}
这里 \texttt{rm} 是机器人上的登录用户名；SSH 服务的安装与登录见第~\ref{sec:deploy-ssh}~节。\texttt{ip route get} 应显示通往机器人的实际网卡与源地址；\texttt{ping} 收到回复说明此刻 ICMP 能往返，但不能证明 SSH 已启动。反过来，防火墙也可能禁 ICMP 而允许 SSH，因此 \texttt{ping} 不通时还要结合实际服务判断。
\end{example}

如果 IP 能连，每次还得问“今天机器人分到几号地址”就有点麻烦了。我们可以给机器人取一个稳定的\textbf{主机名（Hostname）}，例如 \texttt{rm-robot}。但名字不会自动变成地址：直接写 \texttt{ssh rm@rm-robot}，需要本地 DNS、搜索域或 \path{/etc/hosts} 等机制能够解析这个名字。

在常见的小型局域网内，更方便的是多播 DNS（Multicast DNS，mDNS）：用 \texttt{rm-robot.local} 在本地链路查询机器人的地址。Linux 常用 \href{https://avahi.org/}{Avahi} 发布名字，并通过系统解析器或 \href{https://github.com/avahi/nss-mdns}{nss-mdns} 让普通程序也能解析 \texttt{.local}。它要求目标端发布名字、客户端支持解析，且网络放行本地 mDNS 多播；通常不能指望它跨 VLAN、跨路由或穿过隔离网络。

\begin{example}[让 Ubuntu 22.04 机器人能按名字访问]
前提是两端已接入可互访的局域网，机器人已具有不同于队内其他设备的名字。在\textbf{机器人本地终端}配置：
\begin{tbash}
sudo hostnamectl set-hostname rm-robot
sudo apt update
sudo apt install avahi-daemon libnss-mdns
sudo systemctl enable --now avahi-daemon
hostnamectl --static
\end{tbash}
若 \path{/etc/hosts} 的 \texttt{127.0.1.1} 行仍写着旧主机名，用 \texttt{sudoedit /etc/hosts} 同步改为新名，保留原有其他条目。修改名字会影响后续连接，请在本地完成后再重新建立远程会话。

在\textbf{Ubuntu 笔记本}安装同样两个包并启动 Avahi，然后测试：
\begin{tbash}
getent hosts rm-robot.local
ssh rm@rm-robot.local
\end{tbash}
\texttt{getent hosts} 走系统的名字解析路径，出现机器人的地址后才继续 SSH。Windows 或 macOS 笔记本也可以先试 \texttt{ssh rm@rm-robot.local}；能否解析取决于该系统及网络配置，失败时先用已核实的 IP 连接。
\end{example}

如果数字 IP 可用但 \texttt{.local} 不可用，故障就在名字解析这一步附近。先看 \texttt{systemctl status avahi-daemon}，再看 \path{/etc/nsswitch.conf} 的 \texttt{hosts:} 行是否包含可用的 mDNS 解析模块；Ubuntu 常见的是 \texttt{mdns4\_minimal}。不要直接拿网上的一整行覆盖本机配置，因为还可能存在 \texttt{resolve}、\texttt{myhostname} 等系统模块。\texttt{nslookup} 直接询问 DNS 服务器，不能用它的失败单独判定 mDNS 失败；这也是这里选 \texttt{getent} 的原因。

\begin{exercise}\label{ex:deploy-lan}
在自己的网络记录两端地址、前缀及通往机器人的出口网卡，分别用数字 IP 与 \texttt{.local} 名字尝试 SSH。若只有名字失败，给出两项针对名字解析的检查；若都失败，给出两项针对链路、子网或隔离的检查。没有机器人时可用第二台练习电脑，注明替代条件。
\end{exercise}

\section{一根网线直连：不用路由器也能调车}
\label{sec:deploy-wired}

训练场没有可靠的 Wi-Fi，或者临近比赛不便使用无线网络时，先别急着搬出显示器。\textbf{只要机器人系统、网卡和 SSH 还正常，一根网线直连也能使用同样的远程工具}。笔记本没有网口，就加一个 USB 以太网适配器；它提供网卡，不会自动给机器人分配地址。普通笔记本与机器人直连时通常没有 DHCP 服务器，所以我们给两端各设一个固定地址。

下面只演示\textbf{两端都由 NetworkManager 管理的以太网接口}。Ubuntu Desktop 常见这种配置；有些机载 Ubuntu Server 使用 systemd-networkd/Netplan，此时 \texttt{nmcli device status} 可能显示接口未托管，应按现有网络管理方式配置，不能一边保留原管理器一边强行接管。\href{https://networkmanager.pages.freedesktop.org/NetworkManager/NetworkManager/nm-settings-nmcli.html}{NetworkManager 配置参考}列出了地址、默认路由和连接配置的含义。

\begin{note}
切换网卡配置会断开该网卡原来的连接。先在两端本地终端或硬件 KVM 中操作，记录当前连接名，并确认示例子网 \texttt{192.168.77.0/24} 不与现有 LAN、VPN 路由重叠。不要通过唯一的 SSH 会话修改它赖以连接的网卡；也不要在机器人运动时挂着调试线。
\end{note}

\begin{example}[新建独立的直连配置，保留原配置]
先在\textbf{两端}运行以下只读命令，记下以太网接口名和当前连接名：
\begin{tbash}
nmcli device status
nmcli connection show --active
ip route
\end{tbash}
下面假设笔记本网口是 \texttt{enp3s0}，机器人网口是 \texttt{enp2s0}；必须替换为各自实际名称。先插好两端网线，避免启用以太网配置时因没有链路而失败，再在\textbf{笔记本}创建配置并启用：
\begin{tbash}
sudo nmcli connection add type ethernet ifname enp3s0 \
  con-name rm-direct connection.autoconnect no \
  ipv4.method manual ipv4.addresses 192.168.77.1/24 \
  ipv4.never-default yes ipv6.method disabled
sudo nmcli connection up rm-direct
\end{tbash}
\Needspace{9\baselineskip}
再在\textbf{机器人}本地创建它的配置并启用：
\begin{tbash}
sudo nmcli connection add type ethernet ifname enp2s0 \
  con-name rm-direct connection.autoconnect no \
  ipv4.method manual ipv4.addresses 192.168.77.2/24 \
  ipv4.never-default yes ipv6.method disabled
sudo nmcli connection up rm-direct
\end{tbash}
这两段均\textbf{没有设置网关和 DNS}；\texttt{never-default yes} 使这个配置不抢默认路由，笔记本仍可用自己的其他网络上网。\texttt{autoconnect no} 表示下次需手动启用该练习配置；这里暂时禁用该配置的 IPv6，只为让初学实验集中观察 IPv4，并不要求全系统禁用 IPv6。

双方配置启用后，在笔记本验证：
\begin{tbash}
ip route get 192.168.77.2
ping -c 3 192.168.77.2
ssh rm@192.168.77.2
\end{tbash}
结果应显示直连网卡，SSH 应登录机器人；两端没有接入互联网也能完成这些操作。结束时，在两端本地分别执行：
\begin{tbash}
sudo nmcli connection down rm-direct
sudo nmcli connection up "原来的连接名"
\end{tbash}
最后一行的名字换成实验前记录的值；若接口原先没有连接，就无需恢复一个不存在的配置。确认不再需要练习配置后，可用 \texttt{sudo nmcli connection delete rm-direct} 删除\textbf{本次新建的配置}，不要删除原配置。
\end{example}

Windows 笔记本可以在“网络连接”中找到对应的以太网适配器，打开 IPv4 属性，将地址设为 \texttt{192.168.77.1}、掩码设为 \texttt{255.255.255.0}，网关与 DNS 留空；机器人侧仍按上述配置。练习前记录原设置，结束后恢复原来的自动获取或固定配置。注意选择真正插着网线的适配器，别把 Wi-Fi 或 Tailscale 的虚拟网卡改了；具体界面可查\href{https://support.microsoft.com/en-us/windows/experience/connectivity-networking/essential-network-settings-and-tasks-in-windows}{微软的 TCP/IP 设置说明}。

\begin{exercise}\label{ex:deploy-wired}
写一张双端直连配置卡：列出实际接口名、原连接名、两端地址/前缀及恢复命令。在停机的练习设备完成直连，记录路由与 SSH 结果，并说明为什么无需填写网关。如果没有可用网口，只提交配置卡并明确“未实接”，不能用纸面方案替代接线验收。
\end{exercise}

\section{不在同一局域网：用 Tailscale 建立远程维护通道}
\label{sec:deploy-tailscale}

机器人在实验室，你在宿舍；两边都有互联网，却接着不同路由器。这时机器人的 \texttt{192.168.*.*} 私有地址不能直接从外部访问，\texttt{.local} 也通常到不了。\textbf{Tailscale} 可以在获准加入的设备之间建立加密网络，设备组成的网络称为 Tailnet。你仍然运行原来的 SSH 或远程桌面，只是把目标换成 Tailscale 地址或名字，不必把 SSH 端口直接转发到公网。

图~\ref{fig:deploy-network-paths}~把三种情况放在一起：局域网使用现有 LAN 地址；网线直连使用约好的静态地址；跨网络时则在双方联网基础上增加 Tailscale。它们解决“怎么到达目标”，SSH 和远程桌面解决“到达后怎样操作”。

\begin{figure}[htbp]
  \centering
  \includegraphics[width=0.94\textwidth]{deployment/network_paths_revised.png}
  \caption{三种网络入口，最终都连接到机器人上已运行的服务。地址为教学示例；跨网络的 Tailscale 连接可能直达，也可能经过中继。}
  \label{fig:deploy-network-paths}
\end{figure}

\subsubsection{安装、加入同一个 Tailnet，再验证地址}
先和队内负责人确认使用哪个 Tailnet、怎样给成员授权。个人练习可在两端登录同一账户；团队设备应按管理员的成员、标签及访问策略管理，不能把共享登录密码写进教材工程。\href{https://tailscale.com/download}{官方下载页}提供 Windows、macOS 和 Linux 客户端；Windows/macOS 笔记本安装后从托盘或菜单栏登录。

\begin{example}[在 Ubuntu 22.04 机器人安装 Tailscale]
本例采用\href{https://pkgs.tailscale.com/stable/#ubuntu-2204}{官方 Ubuntu 22.04 稳定仓库}，而不是固定某个过时安装包。先在机器人本地或已有 SSH 中确认系统版本：
\begin{tbash}
. /etc/os-release
printf '%s %s\n' "$ID" "$VERSION_CODENAME"
\end{tbash}
下面命令\textbf{只用于输出为 \texttt{ubuntu jammy} 的环境}；其他版本按官方页面选择对应仓库。添加仓库需要管理员权限，会改变系统的软件源：
\begin{tbash}
sudo apt update
sudo apt install curl ca-certificates
curl -fsSL https://pkgs.tailscale.com/stable/ubuntu/jammy.noarmor.gpg \
  -o /tmp/tailscale-jammy-key.gpg
curl -fsSL https://pkgs.tailscale.com/stable/ubuntu/jammy.tailscale-keyring.list \
  -o /tmp/tailscale-jammy.list
\end{tbash}
确认下载成功，查看 \texttt{cat /tmp/tailscale-jammy.list}，其中应指向官方稳定源、\texttt{jammy} 与签名密钥路径；若已配置 Tailscale 软件源，先核对现有设置，不要重复添加。再安装密钥和源文件：
\begin{tbash}
sudo install -d -m 0755 /usr/share/keyrings
sudo install -m 0644 /tmp/tailscale-jammy-key.gpg \
  /usr/share/keyrings/tailscale-archive-keyring.gpg
sudo install -m 0644 /tmp/tailscale-jammy.list \
  /etc/apt/sources.list.d/tailscale.list
sudo apt update
sudo apt install tailscale
sudo systemctl enable --now tailscaled
sudo tailscale up
\end{tbash}
终端会给出登录链接。机载电脑没开浏览器也没关系，在自己的笔记本打开该链接，用计划加入的账户登录，并按管理员要求完成设备批准。笔记本若也是 Ubuntu，同样按其系统版本安装并登录。

在\textbf{机器人}运行：
\begin{tbash}
tailscale status
tailscale ip -4
\end{tbash}
确认设备已登录，记下它实际获得的 Tailscale IPv4 地址。在\textbf{笔记本}运行下面命令；\texttt{100.88.0.20} 仅为示例，必须换成刚才的实际地址：
\begin{tbash}
tailscale status
tailscale ping 100.88.0.20
ssh rm@100.88.0.20
\end{tbash}
\texttt{status} 列出网络中的可见设备和状态；\texttt{tailscale ping} 检查两端 Tailscale 进程的通信路径，可能显示直连，也可能显示中继。它成功仍不等于 SSH 的端口、账户和权限都正常，最后必须用真正的 SSH 登录验证。
\end{example}

\subsubsection{MagicDNS：跨网络也可以用名字}
Tailscale 的 \textbf{MagicDNS} 为设备维护名字。管理员在管理控制台的 DNS 页面确认 MagicDNS 已启用，如图~\ref{fig:deploy-magicdns}；再从设备页面读取机器人的实际设备名及完整域名。它们可能与 Linux 主机名不同，重名时尤其要核对，不能靠猜。

\begin{figure}[htbp]
  \centering
  \includegraphics[width=0.84\textwidth]{deployment/tailscale_magicdns.png}
  \caption{Tailscale 官方文档中的 MagicDNS 启用入口。先看开关状态，再查设备名；已启用的网络无需重复操作。图片来源：\href{https://tailscale.com/docs/features/magicdns}{Tailscale MagicDNS 文档}，界面随版本可能调整。}
  \label{fig:deploy-magicdns}
\end{figure}

\begin{example}[区分两套名字，并检查解析]
假设 Tailscale 控制台显示机器人名为 \texttt{rm-robot}，且客户端接受 Tailnet DNS 配置，可在笔记本运行：
\begin{tbash}
tailscale ping rm-robot
ssh rm@rm-robot
\end{tbash}
若短名与局域网名字冲突，改用控制台显示的完整 MagicDNS 域名。例如 \path{rm-robot.tail-demo.ts.net}，其中 \path{tail-demo} 要换成自己的域名部分，再用实际 Tailscale 数字地址作对照。\path{rm-robot.local} 属于前面的本地 mDNS，后缀为 \path{ts.net} 的完整设备名属于 MagicDNS，二者的解析范围与前提不同。
\end{example}

\subsubsection{连通、授权与赛前检查要分别确认}
本章沿用的是\href{https://tailscale.com/docs/reference/ssh-over-tailscale}{通过 Tailscale 网络访问普通 OpenSSH}：机载电脑上的 \texttt{sshd} 继续负责 SSH 登录认证。另一项名为\href{https://tailscale.com/docs/features/tailscale-ssh}{Tailscale SSH}的功能会接管 Tailnet 上的 SSH 身份与授权，需要专门启用与配置。本章不执行 \texttt{tailscale set --ssh}，免得你在还没理清认证时把两条路线混在一起。

访问还会受到 Tailnet 的 Grants/ACL、设备批准及目标主机防火墙约束。若 \texttt{tailscale ping} 成功而 SSH 超时，核对目标 22 端口、\texttt{sshd}、主机防火墙和 Tailnet 的访问规则；若报 \texttt{Permission denied}，先回到 SSH 账户与密钥。\href{https://tailscale.com/docs/reference/troubleshooting/network-configuration/tcp-connection-two-devices}{官方连通性排错文档}说明了 Tailscale 自身通信与主机 TCP 服务之间的区别。

部署机器人时，还要在控制台检查\href{https://tailscale.com/docs/features/access-control/key-expiry}{设备密钥到期时间}。到期会中断设备连接，重新认证可能使当前远程会话掉线；提前在有本地或有线备用入口时处理。管理员可以对受信任的长期设备关闭密钥到期，但同时需要安排设备丢失后的撤销，不应把关闭到期当成所有设备的必选项。

\begin{conclusion}
局域网名字需要可用的本地解析；\texttt{.local} 通常只在本地链路生效；MagicDNS 依赖 Tailnet 及其 DNS 配置。先用数字地址定位网络问题，再查名字，最后查服务与认证。

Tailscale 适合跨网络维护，不会凭空创造互联网。双方断网、出口封锁或中继链路不可用时，不能保证远程连接；比赛前应实测并准备有线入口。SSH 能连接也不能推出 ROS 2 的自动发现、多播通信或大带宽相机流已可用；这些需要按卷二的通信机制另行配置与验收。
\end{conclusion}

\begin{exercise}\label{ex:deploy-tailscale}
在获准的 Tailnet 内记录机器人实际 Tailscale 地址、设备名、MagicDNS 完整域名与密钥到期状态。把笔记本切到另一条互联网连接，依次验证 \texttt{status}、\texttt{tailscale ping} 和 SSH；说明各自成功能证明什么。没有获准账户或第二条网络时，列出待办与缺少的条件，不能宣称跨网测试已通过。
\end{exercise}

\section{连上以后，把哪一版程序送上车}
\label{sec:deploy-delivery}

远程桌面打开了，终端也连通了，接下来才是部署（Deployment）：把明确版本的代码、依赖、模型和配置放到目标机器，在那里构建或安装，再核对程序是否按约定运行。\textbf{连得上只是维护入口可用，不是视觉系统已经能上场。}

先记录三件事：要部署哪个 Git 提交，目标平台和依赖是什么，如何判断这次部署成功。笔记本常见的 \texttt{x86\_64} 与机载 ARM 的 \texttt{aarch64} 不是同一架构；把前者的可执行文件复制过去，可能直接得到 \texttt{Exec format error}。动态库、GPU 驱动、CUDA、相机 SDK、模型推理后端也各有版本约束。初学时优先在目标机用源码重新构建，交叉编译及板卡工具链留给专项工程。

\subsubsection{用独立目录试跑，再保留能回去的版本}

同步文件可用 \texttt{scp}，基本形式是 \texttt{scp 本地文件 用户@地址:远端目录/}；复制目录加 \texttt{-r}。文件较多、会反复部署时可用 \texttt{rsync}，它能比较并传送变化。通过 SSH 传输时，基本形式为 \texttt{rsync [选项] 源目录/ 主机:目标目录/}。\texttt{-a} 保留常见文件属性，\texttt{-v} 显示过程，\texttt{-n} 只预览。两端都需要安装 \texttt{rsync}；下面使用第~\ref{sec:deploy-ssh}~节的 SSH 别名。

\begin{example}[先同步安全的小项目，再处理真实视觉工程]
在 Ubuntu 笔记本安装 \texttt{rsync}，并在机器人上安装同名软件包。进入本卷目录，确认已有 \path{examples/deployment-demo/}，然后在\textbf{笔记本}执行：
\begin{tbash}
ssh rm-robot 'mkdir -p ~/rm-deploy-demo'
rsync -avn examples/deployment-demo/ rm-robot:rm-deploy-demo/
rsync -av examples/deployment-demo/ rm-robot:rm-deploy-demo/
ssh rm-robot
cd ~/rm-deploy-demo
python3 heartbeat.py --help
\end{tbash}
先看预览列出的文件，再去掉 \texttt{-n} 真正传输。源目录末尾的斜杠表示“把里面的内容放过去”，这里应得到 \path{~/rm-deploy-demo/heartbeat.py}，而不是再套一层 \path{deployment-demo/}。本例故意不加删除目标文件的选项；重复同步已有文件仍会覆盖同名内容，先确认目标是专用练习目录。

此时运行帮助只能证明文件到位和 Python 能启动。下一节还要实际运行、观察日志与停止，再演示服务管理。用虚拟机或 \texttt{localhost} 测试传输，只证明软件流程，不证明机载电脑性能。
\end{example}

真实 C++ 工程也采用这个思路：在目标机建立新的候选目录，从确定的提交取源码；在该目录单独配置、构建和安装。不要把笔记本的 \path{build/} 一起搬过去，CMake 缓存可能记录了笔记本上的编译器和绝对路径。第 3 章的构建命令仍然适用，只是这次在\textbf{机器人 SSH 终端}运行：
\begin{tbash}
uname -m
cmake -S . -B build -DCMAKE_BUILD_TYPE=Release
cmake --build build -j2
cmake --install build --prefix "$HOME/rm-candidate/install"
\end{tbash}
这里假定当前目录就是带有安装规则的候选工程，\path{~/rm-candidate/install} 是本次专用目录。\texttt{-j2} 用两个并行任务，内存紧张的板子可从 \texttt{-j1} 开始；开发机上“线程全开”的习惯搬到小板上，可能先把内存吃光。安装后还需按工程说明传入配置、模型与测试数据，从真实安装位置试跑。

版本不能只靠文件夹名字记。保存源码提交号、未提交修改状态、模型与配置的校验值，以及目标机环境；模型太大不适合直接放 Git 时，也应记录可追溯的存储位置与散列值。修改配置同样会改变行为，不能只记“源码没动”。

\begin{tbash}
git rev-parse HEAD
git status --short
sha256sum config.yaml model.onnx
uname -m
cat /etc/os-release
\end{tbash}
上面最后一组命令在\textbf{真实工程根目录}执行，文件名换成实际的配置与模型。它们不是心跳项目的文件。把输出和启动命令、测试时间、实际结果写入部署记录，便于队友复现。

\begin{exercise}\label{ex:deploy-delivery}
把心跳项目同步到另一台机器的专用目录，确认层级正确，并比较两端 \texttt{heartbeat.py} 的 SHA-256 值。再为自己的 C++ 项目写一份部署记录模板，至少包含版本、架构、依赖、配置、启动方式、验证结果和回退办法。不要把当前可用的比赛目录拿来做覆盖练习。
\end{exercise}

\begin{conclusion}
保留当前已验证版本，在独立目录试跑候选；验收通过后再按队内流程切换。回退要能恢复\textbf{程序、配置与模型的配套组合}，并重新核对启动入口。旧版本放在 Git 提交、发布目录或规范的制品存储中，不靠在源码树里不断增加 \texttt{old2\_final\_reallyfinal} 来维持记忆。
\end{conclusion}

\subsubsection{软件通过以后，还要看真实输入与连续运行}

先用固定录像或心跳程序确认部署流程，再接相机、串口与控制。检查设备权限、相机实际分辨率和帧率、输入输出时间戳、CPU/GPU 占用、内存和温度，以及掉线后的恢复。长时间测试要记录约定时长和实际结束状态；只截一张“进程还在”的图，无法证明整个期间都在正常处理数据。

\begin{example}[给一次板端试跑写清证据]
试跑前，由负责人确定机器人固定方式、允许的动作与停止入口；涉及执行器时，先与电控同学核对安全测试状态。记录候选提交、参数和输入，先跑固定回放核对输出，再在允许状态接入真实设备。试跑过程中保存启动日志、设备断连与重连日志、处理速率及延迟口径、温度和资源采样。结束后写明实际时长、是否异常退出、如何停止，以及是否恢复到当前可用版本。

没有相机或板子时，可以完成文件传输、环境核对与软件回放，把相机输入、功耗温度、真实串口与连续运行列为待验证；不能用“回放通过”替它们打勾。
\end{example}

\begin{exercise}\label{ex:deploy-acceptance}
把“机器人上能跑就行”改写成可检查的验收要求。任选一个视觉功能，分别写出文件与启动、固定输入、真实设备、约定时长运行四类证据，并指出哪两项不能由自己的笔记本回放替代。阈值由实际项目确定，不猜一组好看的数字。
\end{exercise}

文件传输选项与目录末尾斜杠的细节可查 \href{https://download.samba.org/pub/rsync/rsync.1}{rsync 官方手册}；硬件相关依赖与测试则以目标板及对应 SDK 文档为准。

\section{运维工具：保留现场、管好进程与日志}
\label{sec:deploy-ops}

连接机器人以后，我们经常会同时开着程序、看着日志，再改一点参数。SSH 一断，程序还在不在？换个终端启动，怎么就找不到配置了？这时候需要的不只是“再连一次”，还要把进程怎样启动、怎样停止、出了什么问题讲清楚。下面用一个不会访问硬件的心跳程序练习，先把操作跑明白，再迁移到队内工程。

\subsubsection{用 tmux 保留调试终端}

终端复用器（Terminal Multiplexer）tmux 可以让一个远程终端容纳多个窗口，也能把会话暂时放到后台，稍后再接回来。你可以把它想成留在机器人上的一张工作台：笔记本暂时合上了，工作台上的终端还在。程序必须在 tmux 会话里面启动，这个作用才覆盖它；已经在普通 SSH 终端跑起来的进程不会自动搬进去。

\Needspace{11\baselineskip}
\begin{example}[断开以后接回同一个任务]
本卷 \path{examples/deployment-demo} 只有下面三个文件。先把整个目录复制到机器人家目录并命名为 \path{rm-deploy-demo}；确认没有同名旧实验，避免覆盖已有工作。
\begin{lstlisting}[language={}]
rm-deploy-demo/
  heartbeat.py           # 输出运行状态
  rm-heartbeat.service   # 用户服务配置
  README.md              # 复现、停止与清理说明
\end{lstlisting}
Ubuntu 缺少工具时先执行 \texttt{sudo apt install tmux htop}。随后在机器人终端创建名为 \texttt{rm-debug} 的会话并启动程序：
\begin{tbash}
cd ~/rm-deploy-demo
tmux new -s rm-debug
python3 heartbeat.py
\end{tbash}
看到 \texttt{heartbeat=1}、\texttt{heartbeat=2} 持续增加后，先按 \texttt{Ctrl+B}，松开，再按 \texttt{D}，回到原终端。此时可以退出 SSH，重新连接后执行：
\begin{tbash}
tmux ls
tmux attach -t rm-debug
\end{tbash}
如果接回后序号继续增长，而不是又从 1 开始，说明看见的仍是原进程。最后按 \texttt{Ctrl+C} 停止心跳，再输入 \texttt{exit} 结束这个窗口。通用写法是 \texttt{tmux new -s 会话名} 创建，\texttt{tmux attach -t 会话名} 接回。更多窗口和分屏操作见\href{https://github.com/tmux/tmux/wiki/Getting-Started}{tmux 官方入门}。
\end{example}

\begin{note}
tmux 保留的是当前主机上的终端任务。主机关机、系统重启、进程崩溃或系统结束该用户会话，都可能让任务结束。它不能替代服务管理，也不保证退出所有登录后用户进程一定保留；是否结束还受系统登录策略影响。先在自己的目标系统上实际断开、接回验证。
\end{note}

\begin{exercise}\label{ex:deploy-tmux}
在 tmux 中运行心跳，记下 PID 与序号，暂时离开并重新接回，确认是否还是同一进程。随后分别用 Ctrl+C 和 tmux 的离开操作，记录哪一个停止了心跳。再说明重启机器人后，这个会话能不能靠 \texttt{attach} 找回来。
\end{exercise}

\subsubsection{用 systemd 把启动条件写下来}

调试阶段适合手动开终端；需要固定方式运行的任务，可以交给服务管理器 systemd。服务配置把启动命令、运行目录、环境和失败后的动作放在同一个文件里，队友不用再猜“你刚才先敲了哪三条命令”。不过服务通常不会读取交互终端的 \path{.bashrc}，也不会自动激活 Conda 或加载 ROS 环境；这些前提要在启动脚本或服务配置中明确写出来。

先看心跳程序的完整代码。它只用 Python 标准库，打印 PID、版本标记、工作目录和 UTC 时间戳，然后按间隔输出序号；\texttt{flush=True} 让日志及时输出。\texttt{--count} 控制次数，0 表示持续；\texttt{--fail-after} 在第几次输出后以非零状态退出，用来观察失败处理，不是机器人故障检测器。
\lstinputlisting[language=Python,numbers=left,numberstyle={\color{gray}\tiny}]{examples/deployment-demo/heartbeat.py}

\begin{example}[先验证程序，再交给服务]
在机器人 \path{~/rm-deploy-demo} 里先手动运行：
\begin{tbash}
python3 heartbeat.py --count 3 --interval 0.2
python3 heartbeat.py --count 5 --interval 0.2 --fail-after 2
\end{tbash}
第一条应输出三次并正常退出；第二条应在第二次后打印模拟失败，退出状态为 1。每条结束后单独运行 \texttt{echo \$?} 可以查看状态码。参数错误返回非零，持续运行时按 Ctrl+C 正常结束。这个程序没有控制任何设备，适合在笔记本上先练。

接下来使用同目录中的用户服务配置。\texttt{\%h} 是 systemd 展开的用户家目录，\texttt{WorkingDirectory} 决定程序从哪里出发找相对路径；\texttt{ExecStart} 使用解释器的绝对路径，随后给出完整脚本路径。\texttt{Environment} 明确本次运行的版本标记，\texttt{-u} 让 Python 的输出不滞留在缓冲区。
\lstinputlisting[language={},numbers=left,numberstyle={\color{gray}\tiny}]{examples/deployment-demo/rm-heartbeat.service}
确认 \path{~/.config/systemd/user/rm-heartbeat.service} 没有被其他任务占用后，安装并启动：
\begin{tbash}
mkdir -p ~/.config/systemd/user
cp rm-heartbeat.service ~/.config/systemd/user/
systemd-analyze --user verify ~/.config/systemd/user/rm-heartbeat.service
systemctl --user daemon-reload
systemctl --user start rm-heartbeat.service
systemctl --user status rm-heartbeat.service --no-pager
journalctl --user -u rm-heartbeat.service -n 10 --no-pager
\end{tbash}
应看到服务为 \texttt{active (running)}，日志中的版本变成 \texttt{service-v1}，运行目录是用户家目录下的 \path{rm-deploy-demo}。这同时检验了“程序启动”和“启动条件正确”，比只看有没有一个进程更有用。\texttt{verify} 只检查配置，不能代替真正启动。出现“无法连接用户总线”时，先确认是否是正常的 Linux 登录环境、用户管理器是否运行；不要改用 \texttt{sudo systemctl} 把它悄悄变成另一种服务。

修改 unit 后先 \texttt{daemon-reload}，再用 \texttt{systemctl --user restart rm-heartbeat.service} 重启；只改脚本时重启即可。观察结束后执行：
\begin{tbash}
systemctl --user stop rm-heartbeat.service
systemctl --user is-active rm-heartbeat.service
\end{tbash}
应看到 \texttt{inactive}，且日志不再增加心跳。清理这次复制的 unit 与练习文件的方法见项目 README。
\end{example}

\texttt{Restart=on-failure} 表示进程失败后可由管理器重启，\texttt{RestartSec=2s} 留出等待时间；显式 \texttt{stop} 不会因此把任务立即拉回来。本例还限制 30 秒内最多启动 3 次，避免错误配置导致无休止快速重启。你可以在练习副本的 \texttt{ExecStart} 末尾加 \texttt{--fail-after 2}，重载后观察失败与重启记录；触发启动频率限制后，修复原因，再停止并执行 \texttt{systemctl --user reset-failed rm-heartbeat.service}，才重新启动。不要把“它会自动重启”当成故障已经解决。

\begin{note}
这里使用 \texttt{systemctl --user}，只执行 \texttt{start}，没有设置自启动。\texttt{enable} 决定服务是否随相应目标启动；用户服务的生命周期还取决于用户管理器，不能据此保证开机无人登录时运行。真实无人值守部署由队内维护者决定使用系统服务还是用户 lingering，并在目标机器验证启动顺序、设备权限和退出策略。涉及电机或执行机构的程序，应先落实队内安全测试状态，再安排自动启动与失败恢复。
\end{note}

更完整的参数含义可查\href{https://www.freedesktop.org/software/systemd/man/latest/systemd.service.html}{systemd.service}与\href{https://www.freedesktop.org/software/systemd/man/latest/systemd.exec.html}{systemd.exec}。ROS 2 工程还需按卷二的环境加载方式组织启动，不能只把这里的 Python 文件名换成节点名就算部署完成。

\begin{exercise}\label{ex:deploy-service}
在练习副本中修改环境版本标记，重载、重启并从日志核对；再加入 \texttt{--fail-after 2} 观察退出与重启。说明 \texttt{WorkingDirectory}、\texttt{ExecStart}、\texttt{start}、\texttt{enable} 各负责什么，最后恢复配置并停止练习服务。不能运行用户服务的环境，记录实际错误并给出待验证步骤，不把静态检查记成启动成功。
\end{exercise}

\subsubsection{用资源与日志记录定位问题}

程序能启动以后，还要看看它有没有把机载电脑吃满。\texttt{htop} 可以观察进程与 CPU、内存占用；\texttt{df -h} 看文件系统是否快满，\texttt{du -h --max-depth=1 目录} 找出占空间的子目录。它们回答的问题不同：CPU 满了不等于一定需要换模型，磁盘满了也不能先随便删日志。先记录运行场景、进程和时间，再结合算法帧率、延迟与温度看原因；GPU 指标则按实际平台使用厂商工具。

\begin{example}[留下一份能继续排查的记录]
假设心跳已经作为用户服务运行，在机器人终端执行：
\begin{tbash}
systemctl --user show rm-heartbeat.service -p MainPID -p NRestarts
htop
\end{tbash}
第一条给出主进程 PID 和管理器记录的重启次数，在 htop 中按 PID 找到对应进程，观察后按 \texttt{Q} 退出。再在自己的家目录里保留本次信息：
\begin{tbash}
mkdir -p ~/deploy-notes
date -Is > ~/deploy-notes/observed-at.txt
df -h > ~/deploy-notes/disk.txt
journalctl --user -u rm-heartbeat.service -b -n 100 --no-pager > ~/deploy-notes/heartbeat.log
journalctl --user -u rm-heartbeat.service -f
\end{tbash}
\texttt{-b} 选本次系统启动，\texttt{-n 100} 保留最近一百条；\texttt{-f} 持续跟随日志，按 Ctrl+C 只是退出查看，心跳仍由服务管理器运行。队内真实工程还应保存 Git 提交、设备、配置和复现命令。日志能否跨重启保存取决于 journal 的存储配置；需要长期留存时应确认容量和轮转策略，并导出相关片段。分享前检查日志有没有密码、令牌或无关个人信息。
\end{example}

\begin{conclusion}
SSH 是连接入口，tmux 保留调试终端，systemd 管理进程的启动与停止，journalctl 查询其日志。四者可以配合，但“接得上”“进程存在”“日志有心跳”“功能正确”是四个不同判断。排错时沿着连接、启动条件、进程状态和实际输入输出逐层核对，避免只盯着一个绿色 \texttt{active} 放心下车。
\end{conclusion}

\begin{exercise}\label{ex:deploy-ops}
以一次心跳运行建立记录：写出观察时间、运行版本、PID、重启次数、磁盘使用和最近日志。然后只退出 \texttt{journalctl -f}，检查进程是否还在；最后真正停止服务并确认结果。如果真实视觉程序显示 \texttt{active} 却没有新图像，指出还应核对哪两类输入输出证据。
\end{exercise}

\section{赛前与故障救援：让笔记本变成硬件控制台}
\label{sec:deploy-kvm}

快要上场时，无线连接可能受到现场条件或赛事规定限制。看起来又得抱着便携屏和键盘蹲回地上了？先别急：如果机器人系统与网络服务正常，第~\ref{sec:deploy-wired}~节的\textbf{有线 SSH} 就能解决不少问题。若连网卡都没配好、系统还停在启动界面，或你需要操作 BIOS，则可以准备\textbf{键盘、视频、鼠标（Keyboard, Video, Mouse，KVM）}工具。

这里说的是硬件 KVM：它接收目标电脑的视频输出，并模拟 USB 键鼠，把这些信号送到你的笔记本。它不要求机器人先装好远程桌面服务。和前面的 SSH、RDP 相比，这种独立于目标系统网络服务的访问通常称为\textbf{带外管理（out-of-band management）}。

图~\ref{fig:deploy-kvm-wiring}~先把两种接法放在一起。顺着箭头看：机器人视频送到 KVM，键鼠输入则从 KVM 送回机器人；笔记本接的是 KVM 的控制端，不能省掉机器人侧的线。
\begin{figure}[htbp]
  \centering
  \includegraphics[width=0.95\textwidth]{deployment/kvm_wiring_revised.png}
  \caption{USB KVM 与 IP KVM 的接线关系示意。“笔记本侧一根线”与“机器人侧多条连接”同时成立；供电方式须按具体型号核对。}
  \label{fig:deploy-kvm-wiring}
\end{figure}

\subsubsection{USB KVM：一根线接笔记本，另一侧接机器人}

以 \textbf{Openterface Mini-KVM} 为例，图~\ref{fig:deploy-usb-kvm} 展示设备与配套线缆。笔记本安装厂商的控制端应用，通过一根\emph{支持数据传输的} USB 线接到 KVM 的 Host 端，就可以用自己的屏幕和键鼠操作机器人。机器人端则要把 HDMI 视频输出接到 KVM 的 Input，把机器人 USB 口接到 KVM 的 Target，使其收到模拟键鼠输入。机器人的正常电源仍要保留；这台 Mini-KVM 自身可通过 USB 供电，不代表笔记本也在为机器人供电。

\begin{figure}[htbp]
  \centering
  \includegraphics[width=0.72\textwidth]{deployment/openterface-mini-kvm.jpg}
  \caption{Openterface Mini-KVM 及配套线缆官方照片。橙色长线连接笔记本，短 USB 线与 HDMI 线连接目标电脑；“一根线”描述的是笔记本一侧。}
  \label{fig:deploy-usb-kvm}
\end{figure}

\begin{example}[不依赖网络，操作机载电脑]
先按\href{https://docs.openterface.com/app/kvm/}{官方应用下载页}在笔记本安装对应 Windows、macOS 或 Linux 的 Openterface 控制端。Linux 应选择匹配架构的包，并按其安装说明处理 USB/视频设备权限；机器人端不用安装这款应用。
\begin{enumerate}
  \item 让机器人处于安全测试状态，确认其显示输出接口及所需适配器。在连接处看清 \texttt{Host}、\texttt{Target}、\texttt{Input} 标识，避免只凭 USB-C 外形判断用途。
  \item 用数据 USB 线连接笔记本与 Host；用另一条 USB 线连接机器人与 Target；再把机器人 HDMI 输出接到 KVM 的 HDMI Input。核对两台电脑的供电，不向不明电源接口送电。
  \item 打开笔记本控制端，选择检测到的 KVM 视频设备与控制设备。看到机器人画面后，在应用内点击或输入，打开终端运行 \texttt{hostname}，确认图像和键鼠两条链路都通了。
  \item 在事先约定的维护窗口，可重启\emph{测试机}，观察启动画面并按该主板说明进入 BIOS/UEFI；只查看，不改启动项。由此检查该板在开机阶段能否识别 KVM，不能拿运行中桌面的成功替代这一步。
\end{enumerate}
看得见却点不动时，优先核对 Target USB、数据线与控制权限；完全没有画面时，检查 HDMI 输入、分辨率与输出模式。两者都失败才一起查 Host 端连接和应用设备选择。具体接线对照\href{https://docs.openterface.com/products/minikvm/how-to-connect/}{官方接线指南}。
\end{example}

\begin{note}
多数笔记本的 HDMI 口是\emph{输出}，拿一根 HDMI 线把两台电脑相连，笔记本通常不会变成显示器。普通 HDMI 采集卡可以把画面送入笔记本，但它没有自动提供 USB 键鼠模拟功能；完整控制还需要相应输入链路。USB-C 同样只是接口形状，普通两台电脑之间也不能靠随便一根 USB 线直接互控。
\end{note}

\subsubsection{IP KVM：通过网线进入浏览器控制台}

另一类是 \textbf{IP KVM（KVM over IP）}，例如 \textbf{PiKVM}。图~\ref{fig:deploy-ip-kvm} 中可看到它自己的以太网口：HDMI 和 USB 仍连接机器人，笔记本则通过网络访问\emph{KVM 本身}的网页。机器人自己的网卡就算配置错了，只要 KVM 的网络、供电与视频输入正常，仍有机会从网页控制台进去修复。

\begin{figure}[htbp]
  \centering
  \includegraphics[width=0.78\textwidth]{deployment/pikvm-v4.jpg}
  \caption{PiKVM V4 Mini（左）与 Plus（右）官方照片。网线连接 KVM 的网络端口；其他 RJ45 外形接口可能承担串口或 ATX 功能，必须按标识区分。}
  \label{fig:deploy-ip-kvm}
\end{figure}

\begin{example}[在有线局域网里使用 PiKVM]
以已装好系统的 PiKVM V4 为例，按型号手册接好机器人 HDMI 输出和 USB 键鼠模拟口，并使用规定的独立电源。用网线将 PiKVM 与笔记本接入同一个带 DHCP 的有线调试网络，从 PiKVM 屏幕或路由器租约表读取其地址，例如 \texttt{192.168.50.30}。

在笔记本浏览器打开 \texttt{https://192.168.50.30/}。原厂默认配置的网页账户为 \texttt{admin}、密码为 \texttt{admin}，双因素认证码默认未启用；已配置过的设备应向管理者取得实际凭据，不能靠反复尝试默认密码接管未知设备。自签名证书提示出现时，先通过本地屏幕、租约与接线核对设备身份。

\Needspace{11\baselineskip}
首次登录后先点 Terminal，进入的是 \emph{PiKVM 自己的管理终端}。运行 \texttt{su -}，原厂系统管理员 \texttt{root} 的默认密码为 \texttt{root}，随后逐行运行：
\begin{tbash}
rw
passwd root
kvmd-htpasswd set admin
ro
exit
\end{tbash}
\texttt{rw}/\texttt{ro} 临时切换 PiKVM 的系统分区可写/只读状态；中间两条命令分别交互式设置系统管理员和网页管理员的新密码。退出终端，重新登录网页验证新密码，然后点击 KVM 进入机器人的控制台，打开目标机终端核对 \texttt{hostname}。网页的 KVM 画面是机器人，PiKVM 的 Web Terminal 是 KVM 管理系统，两者不要弄反。接线、出厂凭据与首次配置见\href{https://docs.pikvm.org/v4/}{PiKVM V4 官方快速指南}。

如果想把笔记本用一根网线\emph{直接}接到 PiKVM，而不经过路由器，应先按对应固件说明为它与笔记本配置同网段静态地址；两台设备都在等待 DHCP 时，插线不会凭空出现 DHCP 服务器。先在有 DHCP 的网络完成首次配置，再验证脱网直连，赛前就不用临场猜地址。
\end{example}

IP KVM 的“一根网线”也只描述笔记本侧的访问链路，KVM 到机器人仍需视频、USB 与相应供电。某些型号还有 ATX 开关机控制，但必须另接受支持的主板接口；有 KVM 不等于自动获得机器人总电源控制。板子没有显示输出、供电损坏，或系统卡死后根本不响应 USB，KVM 也不能把它凭空救活。确认具体主板的输出、开机键盘识别与 KVM 固件兼容性，比背产品参数更有用。

\begin{exercise}\label{ex:deploy-kvm}
给“系统正常但无线不能用”“机器人 IP 配错”“停在 BIOS”“目标机完全断电”四种场景分别选择有线 SSH、USB KVM 或 IP KVM，并说明成立条件。再画一张 KVM 接线图，标清笔记本侧和机器人侧各要几条线、谁提供电源；有设备时现场演示“有画面”和“能输入”两项独立验证。
\end{exercise}

\begin{conclusion}
系统和网络服务正常，先用有线 SSH；需要窗口且桌面正常，用 RDP 或 RustDesk；需要查看启动过程、修改网络或绕过失效的桌面服务，考虑硬件 KVM。USB KVM 让笔记本通过 USB 充当控制台，IP KVM 让浏览器通过网络访问控制台。它们都需要目标机具备可用的视频与输入接口；“一根线”不取消机器人侧的接线与供电。
\end{conclusion}

上场前把线缆、转接头、备用地址、登录方式和恢复步骤装进调试包，在本机演练一次，再按当届赛事规则及现场裁判要求确认哪些连接允许使用、何时必须拔除。不能因为改成有线就默认比赛全过程都允许连接；正式启动前还要清点外接线缆，防止运动时拉扯接口。


\section{本章小结}
\begin{conclusion}
代码在自己的电脑上运行，是实机部署的起点。先为机载电脑准备系统、依赖和维护入口，再把确定版本的代码、配置与模型送过去，在目标环境构建、试跑和记录结果。

便携屏与键鼠适合初次安装和现场救援；日常命令操作用 SSH，需要桌面再用 RDP 或其他远程桌面。网络方面，先确认链路和地址，再处理名字：局域网可用 DNS 或 \texttt{.local}，网线直连要配同一子网的地址，跨网络可用 Tailscale 的私有地址和 MagicDNS。名称、网络可达、服务与身份验证是不同层次，逐层排查比反复重装软件更有效。

临时保留终端用 tmux，长期管理程序用 systemd，靠日志与资源记录判断运行状态；无无线环境可用有线网络，操作系统网络失效时再考虑带外 KVM。连接工具给你操作入口，最终交付仍要经过真实输入、异常恢复和约定时长的验证，并保留可以回退的配套版本。
\end{conclusion}

\begin{problemset}
  \item\label{prob:v1-deploy-1} （基础）笔记本上的程序能运行，复制到 ARM 机载电脑后却报 \texttt{Exec format error}，还可能遇到哪些环境差异？首次安装系统时为什么仍值得准备便携屏与键鼠？
  \item\label{prob:v1-deploy-2} （基础）某人把 \texttt{ssh rm-robot} 中的别名当成真实主机名，又把机器人主机指纹当成自己的公钥。分别纠正这两处混淆，并说明 \texttt{Permission denied} 与连接超时通常应从哪层开始查。
  \item\label{prob:v1-deploy-3} （基础）比较 VS Code Remote -- SSH、GNOME RDP、RustDesk 与 KVM 各自提供的操作入口。Ubuntu 22.04 未登录图形桌面时，为什么不能直接断言其桌面共享一定可用？
  \item\label{prob:v1-deploy-4} （基础）同一路由器下，IP 可以连接而 \texttt{rm-robot.local} 不可解析；换到校园另一网段后，\texttt{.local} 又失效。分别说明应检查什么；裸主机名、\texttt{.local} 与 MagicDNS 是否是同一机制？
  \item\label{prob:v1-deploy-5} （基础）笔记本直连机器人网口，双方都没有 DHCP 服务。请设计不与现有网络冲突的地址与掩码，说明是否需要默认网关、DNS，并写出如何恢复原连接。
  \item\label{prob:v1-deploy-6} （提高）双方均安装 Tailscale，\texttt{tailscale ping} 成功却无法 SSH 登录。给出排查顺序；若机载电脑彻底失去互联网，远方队友能否一定靠 Tailscale 救回来？为什么不能直接用它保证 ROS 2 自动发现？
  \item\label{prob:v1-deploy-7} （基础）解释 \texttt{rsync -avn src/ host:dst/} 的预览与目录含义。部署记录为什么应包含模型、配置与未提交改动？回退时只切换源码提交可能遗漏什么？
  \item\label{prob:v1-deploy-8} （提高）比较 SSH 普通前台进程、tmux 会话与 systemd 服务对断线、崩溃、重启的处理。用户服务文件保存好了却未运行，或打开终端能跑、服务中找不到依赖，分别应查哪些条件？
  \item\label{prob:v1-deploy-9} （基础）“USB KVM 只要一根线”“IP-KVM 插上网线就能操控机器人”“采集卡可以代替 KVM”分别缺少什么前提？画出笔记本、KVM、机器人之间的视频、键鼠与网络连接。
  \item\label{prob:v1-deploy-10} （实践自测）设计一次赛前演练：无线不可用，候选版本启动后没有新日志。给出有线接入、分层定位、恢复可用版本与退出现场连接的方案，写清应留下的记录。可在练习机完成软件部分，涉及比赛与机器人部分列明实际验证条件。
\end{problemset}

\section{本章验收}
下面检查实际操作能力，允许查文档。涉及机器人动作时由负责人确认测试状态与停止入口；没有对应设备的项目保留为待验收，不能拿模拟截图顶替。
\begin{enumerate}
  \item 在两台设备间核验主机指纹并用自己的密钥 SSH 登录，现场指出命令在哪台机器执行；通过别名或可解析的主机名连接，解释所用解析机制。
  \item 在专用连接配置下完成一次网线直连，暂时停用无线连接后仍能 SSH；演示查看实际路由与地址，并恢复原网络配置。跨网条件可用时，再演示 Tailscale 的在线状态、名称和目标端口验证。
  \item 演示一种适合实际系统的远程桌面，指出它依赖的服务、登录状态和网络条件。若具备 USB KVM 或 IP-KVM，再在操作系统网络不可用时演示画面和键鼠控制，说明各根线与供电作用。
  \item 从专用目录部署心跳项目，用 tmux 保留会话，再用 systemd 启停并查询日志。由学长在副本设置一个路径、工作目录或退出故障，依据实际证据定位，恢复后停止练习服务。
  \item 给出一份可以交给队友的部署记录：版本、环境、模型配置、启动与停止、实际验证结果、未验证项和回退入口。真实工程还须另行提交设备输入、异常恢复与约定时长的测试记录。
\end{enumerate}
\endgroup
